% GENERATED FILE. Edit canonical lessons, then run npm run generate:books.
% canonical-source-hash: fnv1a64:337b1142f6e54e58
% canonical-lessons: LA-C321-minae, LA-C321-minitor, LA-C321-terreo, LA-C321-tumultus, LA-C321-clamor

\chapter{Minae et tumultus}
\label{ch:la-minae-et-tumultus}
\hlchaptermodality{321}

\begin{chapteropening}
\textbf{By the end of this chapter:} \emph{I can tell of threats and uproar: threats, to threaten, to frighten, an uproar and shouting.}

Complete the last lesson of chapter 321: I can tell of threats and uproar: threats, to threaten, to frighten, an uproar and shouting.
\end{chapteropening}

\section[minae, minārum]{minae, minārum — threats}

\label{lesson:LA-C321-minae}



\begin{quote}
Before the new one: say the Latin for frightened, then the Latin for to panic.
\end{quote}

\subsection*{You'll want to know: minae, minārum}
\textbf{minae, minārum} — \textquotedblleft{}threats\textquotedblright{}.

\begin{quote}
Hostēs minās contrā urbem iactābant.
\end{quote}

\begin{quote}
\emph{The enemy were hurling threats against the city.}
\end{quote}

Feminine, plural. First the jutting points of a wall, then threats that hang over someone.

\begin{cousinweb}
\textbf{minae} first meant \textquotedblleft{}the jutting points of a wall\textquotedblright{} (Lewis \& Short); EtymDB traces it to Proto-Italic \emph{*menā}, \textquotedblleft{}a projection\textquotedblright{}. English \emph{menace} comes from its family.
\end{cousinweb}

\begin{grammarlens}[title={What you've built}]
One more word for telling, explaining and describing at B1.
\end{grammarlens}

\subsection*{Your turn}
\begin{itemize}
\raggedright
  \item \emph{Say it:} \emph{minae, minārum}
  \item \emph{Say it:} \emph{minae, minārum}, once more
  \item \emph{Say it:} a sentence of your own with \emph{minae, minārum}
  \item \emph{Recall:} say the Latin for frightened, then the Latin for to panic, then say \emph{minae, minārum} again
\end{itemize}

\begin{tcolorbox}[breakable,colback=teal!4,colframe=teal!35!black,title={Before you move on}]
What does minae, minārum mean? (\textquotedblleft{}Threats\textquotedblright{}.) Say it once more.
\end{tcolorbox}
\section[minitor, minitārī]{minitor, minitārī — to threaten}

\label{lesson:LA-C321-minitor}



\begin{quote}
Before the new one: say the Latin for to panic, then the Latin for threats.
\end{quote}

\subsection*{You'll want to know: minitor, minitārī}
\textbf{minitor, minitārī} — \textquotedblleft{}to threaten\textquotedblright{}.

\begin{quote}
Latrō mihi gladiō minitābātur.
\end{quote}

\begin{quote}
\emph{The bandit was threatening me with a sword.}
\end{quote}

\emph{minitor, minitārī, minitātus sum}, a deponent, with the dative of the person threatened.

\begin{cousinweb}
\textbf{minitor} is a stronger form of \emph{minor}, \textquotedblleft{}to threaten\textquotedblright{}, which is built on \emph{minae}, \textquotedblleft{}threats\textquotedblright{} (Lewis \& Short).
\end{cousinweb}

\begin{grammarlens}[title={What you've built}]
One more word for telling, explaining and describing at B1.
\end{grammarlens}

\subsection*{Your turn}
\begin{itemize}
\raggedright
  \item \emph{Say it:} \emph{minitor, minitārī}
  \item \emph{Say it:} \emph{minitor, minitārī}, once more
  \item \emph{Say it:} a sentence of your own with \emph{minitor, minitārī}
  \item \emph{Recall:} say the Latin for to panic, then the Latin for threats, then say \emph{minitor, minitārī} again
\end{itemize}

\begin{tcolorbox}[breakable,colback=teal!4,colframe=teal!35!black,title={Before you move on}]
What does minitor, minitārī mean? (\textquotedblleft{}To threaten\textquotedblright{}.) Say it once more.
\end{tcolorbox}
\section[terreō, terrēre]{terreō, terrēre — to frighten}

\label{lesson:LA-C321-terreo}



\begin{quote}
Before the new one: say the Latin for threats, then the Latin for to threaten.
\end{quote}

\subsection*{You'll want to know: terreō, terrēre}
\textbf{terreō, terrēre} — \textquotedblleft{}to frighten\textquotedblright{}.

\begin{quote}
Tonitrus īnfantem terruit.
\end{quote}

\begin{quote}
\emph{The thunder frightened the baby.}
\end{quote}

\emph{terreō, terrēre, terruī, territum}. English \emph{terrify} and \emph{terror} belong to this family.

\begin{cousinweb}
\textbf{terreō} is inherited from a Proto-Indo-European root meaning \textquotedblleft{}to tremble, to frighten\textquotedblright{} (EtymDB). English \emph{terrify} and \emph{terror} come from its family.
\end{cousinweb}

\begin{grammarlens}[title={What you've built}]
One more word for telling, explaining and describing at B1.
\end{grammarlens}

\subsection*{Your turn}
\begin{itemize}
\raggedright
  \item \emph{Say it:} \emph{terreō, terrēre}
  \item \emph{Say it:} \emph{terreō, terrēre}, once more
  \item \emph{Say it:} a sentence of your own with \emph{terreō, terrēre}
  \item \emph{Recall:} say the Latin for threats, then the Latin for to threaten, then say \emph{terreō, terrēre} again
\end{itemize}

\begin{tcolorbox}[breakable,colback=teal!4,colframe=teal!35!black,title={Before you move on}]
What does terreō, terrēre mean? (\textquotedblleft{}To frighten\textquotedblright{}.) Say it once more.
\end{tcolorbox}
\section[tumultus, tumultūs]{tumultus, tumultūs — an uproar, a commotion}

\label{lesson:LA-C321-tumultus}



\begin{quote}
Before the new one: say the Latin for to threaten, then the Latin for to frighten.
\end{quote}

\subsection*{You'll want to know: tumultus, tumultūs}
\textbf{tumultus, tumultūs} — \textquotedblleft{}an uproar, a commotion\textquotedblright{}.

\begin{quote}
Magnus tumultus in forō erat.
\end{quote}

\begin{quote}
\emph{There was a great commotion in the forum.}
\end{quote}

Masculine, 4th declension. Also an emergency war close to home. English \emph{tumult}.

\begin{cousinweb}
\textbf{tumultus} is perhaps related to \emph{tumeō}, \textquotedblleft{}to swell\textquotedblright{}; its origin is not settled (Lewis \& Short). English \emph{tumult} comes from it through French.
\end{cousinweb}

\begin{grammarlens}[title={What you've built}]
One more word for telling, explaining and describing at B1.
\end{grammarlens}

\subsection*{Your turn}
\begin{itemize}
\raggedright
  \item \emph{Say it:} \emph{tumultus, tumultūs}
  \item \emph{Say it:} \emph{tumultus, tumultūs}, once more
  \item \emph{Say it:} a sentence of your own with \emph{tumultus, tumultūs}
  \item \emph{Recall:} say the Latin for to threaten, then the Latin for to frighten, then say \emph{tumultus, tumultūs} again
\end{itemize}

\begin{tcolorbox}[breakable,colback=teal!4,colframe=teal!35!black,title={Before you move on}]
What does tumultus, tumultūs mean? (\textquotedblleft{}An uproar, a commotion\textquotedblright{}.) Say it once more.
\end{tcolorbox}
\section[clāmor, clāmōris]{clāmor, clāmōris — a shout, shouting}

\label{lesson:LA-C321-clamor}



\begin{quote}
Before the new one: say the Latin for to frighten, then the Latin for an uproar.
\end{quote}

\subsection*{You'll want to know: clāmor, clāmōris}
\textbf{clāmor, clāmōris} — \textquotedblleft{}a shout, shouting\textquotedblright{}.

\begin{quote}
Subitō clāmor ingēns in viā ortus est.
\end{quote}

\begin{quote}
\emph{Suddenly a huge shout went up in the street.}
\end{quote}

Masculine, 3rd declension, from \emph{clāmō} (taught). \emph{Clāmōrem tollere}: to raise a shout. English \emph{clamour}.

\begin{cousinweb}
\textbf{clāmor} is built on \emph{clāmō}, \textquotedblleft{}to shout\textquotedblright{} (EtymDB). English \emph{clamour} comes from it through French.
\end{cousinweb}

\begin{grammarlens}[title={What you've built}]
One more word for telling, explaining and describing at B1.
\end{grammarlens}

\subsection*{Your turn}
\begin{itemize}
\raggedright
  \item \emph{Say it:} \emph{clāmor, clāmōris}
  \item \emph{Say it:} \emph{clāmor, clāmōris}, once more
  \item \emph{Say it:} a sentence of your own with \emph{clāmor, clāmōris}
  \item \emph{Recall:} say the Latin for to frighten, then the Latin for an uproar, then say \emph{clāmor, clāmōris} again
\end{itemize}

\begin{tcolorbox}[breakable,colback=teal!4,colframe=teal!35!black,title={Before you move on}]
What does clāmor, clāmōris mean? (\textquotedblleft{}A shout, shouting\textquotedblright{}.) Say it once more.
\end{tcolorbox}
